\documentclass[12pt,a4paper]{article}

% ─── Packages ────────────────────────────────────────────────────────────────
\usepackage[utf8]{inputenc}
\usepackage[T5]{fontenc}
\usepackage[vietnamese,english]{babel}
\usepackage{amsmath, amssymb, mathtools}
\usepackage{graphicx}
\usepackage{float}
\usepackage{booktabs}
\usepackage{array}
\usepackage{xcolor}
\usepackage{listings}
\usepackage{hyperref}
\usepackage{geometry}
\usepackage{caption}
\usepackage{subcaption}
\usepackage{enumitem}
\usepackage{titlesec}
\usepackage{fancyhdr}
\usepackage{parskip}

\geometry{margin=2.5cm}

% ─── Code style ──────────────────────────────────────────────────────────────
\lstdefinestyle{python}{
  language=Python,
  basicstyle=\ttfamily\small,
  keywordstyle=\color{blue!70!black}\bfseries,
  commentstyle=\color{green!50!black}\itshape,
  stringstyle=\color{red!60!black},
  numberstyle=\tiny\color{gray},
  numbers=left,
  stepnumber=1,
  frame=single,
  breaklines=true,
  tabsize=4,
  showstringspaces=false,
  backgroundcolor=\color{gray!8},
}
\lstset{style=python}

% ─── Header / Footer ─────────────────────────────────────────────────────────
\pagestyle{fancy}
\fancyhf{}
\rhead{MAT3562E -- Computer Vision Lab 02}
\lhead{Bui Quang Chien -- 23001837}
\cfoot{\thepage}

% ─── Hyperref ────────────────────────────────────────────────────────────────
\hypersetup{
  colorlinks=true,
  linkcolor=blue!60!black,
  urlcolor=blue,
}

% ─── Title ───────────────────────────────────────────────────────────────────
\title{%
  \textbf{LAB 02 REPORT}\\[0.4em]
  \large Arithmetic \& Logical Operations \\
  and 2D Convolution \\[0.8em]
  \normalsize MAT3562E -- Computer Vision\\
  Academic Year 2025--2026
}
\author{
  Bui Quang Chien\\
  MSSV: 23001837
}
\date{March 2026}

% ─────────────────────────────────────────────────────────────────────────────
\begin{document}
\maketitle
\tableofcontents
\newpage

% =============================================================================
\section{Overview and Learning Objectives}
% =============================================================================

This lab has two parts:
\begin{itemize}[leftmargin=*]
  \item \textbf{Part 1 -- Pixel-wise Operations:}
        arithmetic operations (gamma correction, add, subtract, multiply, divide),
        logical operations (AND, OR, XOR), and Otsu automatic thresholding.
  \item \textbf{Part 2 -- Neighborhood Operations:}
        2D convolution with VALID and SAME padding modes, using both NumPy and
        OpenCV, with a library of standard kernels (blur, sharpen, edge detection).
\end{itemize}

Each method is implemented in two ways: (A) NumPy-only and (B) OpenCV, for
comparison and validation.

% =============================================================================
\section{Environment and Common Utilities}
% =============================================================================

All code runs in Python~3 with the following packages:
\texttt{numpy}, \texttt{opencv-python}, \texttt{matplotlib}.

The following utility functions are shared by all tasks:

\begin{lstlisting}
import numpy as np, cv2, matplotlib.pyplot as plt

def read_gray_cv2(path):
    img = cv2.imread(path, cv2.IMREAD_GRAYSCALE)
    if img is None: raise FileNotFoundError(path)
    return img                   # uint8 [H,W] in [0..255]

def to_float01(img_u8):
    return img_u8.astype(np.float32) / 255.0    # [0..1]

def to_uint8(img_f):
    return (np.clip(img_f, 0, 1) * 255 + 0.5).astype(np.uint8)

def normalise_for_display(arr):   # any float -> uint8 [0..255]
    mn, mx = arr.min(), arr.max()
    if mx == mn: return np.zeros_like(arr, dtype=np.uint8)
    return ((arr - mn) / (mx - mn) * 255).astype(np.uint8)
\end{lstlisting}

% =============================================================================
\section{Part 1 -- Arithmetic \& Logical Operations}
% =============================================================================

\subsection{Conceptual Motivation}
All operations in Part~1 are \emph{point operations}:
\[
  \text{output}(x,y) = f\!\bigl(\text{input}(x,y)\bigr)
\]
Each pixel is processed independently of its neighbours.

% ─────────────────────────────────────────────────────────────────────────────
\subsection{Gamma Correction (Power-law Transform)}
% ─────────────────────────────────────────────────────────────────────────────

\subsubsection*{Theory}
\begin{equation}
  s = r^{\gamma}, \qquad r \in [0,1]
\end{equation}
\begin{itemize}[leftmargin=*]
  \item $\gamma < 1$: expands dark tones (brightens shadows), compresses highlights.
  \item $\gamma > 1$: compresses shadows, expands highlights (overall darker).
  \item $\gamma = 1$: identity.
\end{itemize}

\subsubsection*{Implementation}

\paragraph{NumPy:}
\begin{lstlisting}
def gamma_correction_np(img_u8, gamma):
    r = to_float01(img_u8)
    s = np.power(r, gamma)     # broadcast, element-wise
    return to_uint8(s)
\end{lstlisting}

\paragraph{OpenCV (LUT-based):}
\begin{lstlisting}
def gamma_correction_cv2(img_u8, gamma):
    table = np.array([((i/255.0)**gamma)*255 for i in range(256)],
                     dtype=np.uint8)
    return cv2.LUT(img_u8, table)   # faster: precomputed lookup table
\end{lstlisting}

\subsubsection*{Mini-Task Results}

\textbf{Image used:} \texttt{2.png} (daytime outdoor/sunrise scene).
\textbf{Parameters tested:} $\gamma \in \{0.5,\; 0.8,\; 1.2,\; 2.0\}$.

\begin{table}[H]
  \centering
  \begin{tabular}{@{}lp{9cm}@{}}
    \toprule
    $\gamma$ & Visual effect \\
    \midrule
    0.5 & Significantly brighter; dark shadow regions become distinguishable; highlights saturate slightly. \\
    0.8 & Mildly brighter; subtle brightening of mid-tones. \\
    1.2 & Mildly darker; highlights compressed. \\
    2.0 & Noticeably darker; shadows become very dense; contrast in bright regions increases. \\
    \bottomrule
  \end{tabular}
  \caption{Gamma correction results on the test image.}
\end{table}

\textbf{Regions that change most:} Pixels in the mid-to-low intensity range change
most drastically for $\gamma = 0.5$ (shadow brightening) and $\gamma = 2.0$
(shadow darkening). The very darkest ($r \approx 0$) and very brightest
($r \approx 1$) pixels are fixed points of $s = r^{\gamma}$.

% ─────────────────────────────────────────────────────────────────────────────
\subsection{Arithmetic Operations}
% ─────────────────────────────────────────────────────────────────────────────

\subsubsection*{Formulas}
\[
  S_{\text{add}} = \alpha A + \beta B, \quad
  S_{\text{sub}} = \alpha A - \beta B, \quad
  S_{\text{mul}} = A \odot B, \quad
  S_{\text{div}} = \frac{A}{B + \varepsilon}
\]

\subsubsection*{NumPy vs OpenCV}

\begin{table}[H]
  \centering
  \begin{tabular}{@{}lll@{}}
    \toprule
    Operation & NumPy (float+clip) & OpenCV (saturating uint8) \\
    \midrule
    Addition    & Soft saturation via clip & Hard saturation at 255 \\
    Subtraction & Soft floor via clip at 0 & Hard floor at 0 \\
    Multiply    & Always $\le 1$ in float & Can overflow; use \texttt{scale} \\
    Division    & May exceed 1; clipped    & Use \texttt{scale} to keep in range \\
    \bottomrule
  \end{tabular}
  \caption{NumPy vs OpenCV arithmetic behaviour.}
\end{table}

\subsubsection*{Guiding Questions}
\begin{enumerate}
  \item \textbf{Why is subtraction useful for change detection?}
        If two aligned images $A$ (background) and $B$ (background + object) are taken
        under identical conditions, $|A - B|$ is near-zero everywhere except where something
        changed.
  \item \textbf{Why can division produce large values?}
        When $B \approx 0$ (dark pixel), $A/(B+\varepsilon)$ can be very large even for
        small $\varepsilon$.
  \item \textbf{When does clipping hide information?}
        When the result meaningfully exceeds [0,255], e.g.\ in multi-image addition or
        division. Min-max normalisation should be used instead.
\end{enumerate}

% ─────────────────────────────────────────────────────────────────────────────
\subsection{Logical Operations (AND / OR / XOR)}
% ─────────────────────────────────────────────────────────────────────────────

\subsubsection*{Workflow}
\begin{enumerate}
  \item Binarise via manual threshold or Otsu.
  \item Apply AND / OR / XOR on binary masks.
  \item Optionally AND the mask with the grayscale image to extract ROI.
\end{enumerate}

\begin{table}[H]
  \centering
  \begin{tabular}{@{}lll@{}}
    \toprule
    Operator & Condition for white pixel & Use case \\
    \midrule
    AND & Both masks are white & Intersection; conservative (fewest FP) \\
    OR  & At least one mask is white & Union; liberal (fewest FN) \\
    XOR & Exactly one mask is white & Disagreement / boundary pixels \\
    \bottomrule
  \end{tabular}
  \caption{Logical operator semantics on binary masks.}
\end{table}

% ─────────────────────────────────────────────────────────────────────────────
\subsection{Otsu Automatic Thresholding}
% ─────────────────────────────────────────────────────────────────────────────

\subsubsection*{Theory}
Otsu maximises the \emph{between-class variance}:
\begin{equation}
  \sigma_b^2(t) = \omega_0(t)\,\omega_1(t)\,\bigl(\mu_0(t) - \mu_1(t)\bigr)^2
\end{equation}
where $\omega_0, \omega_1$ are class probabilities and $\mu_0, \mu_1$ are class means.
Works best for bimodal histograms.

\subsubsection*{NumPy Implementation (from scratch)}
\begin{lstlisting}
def otsu_threshold_np(img_u8):
    hist = np.bincount(img_u8.ravel(), minlength=256).astype(np.float64)
    p    = hist / hist.sum()
    omega = np.cumsum(p)               # cumulative class probability
    mu    = np.cumsum(p * np.arange(256))
    mu_T  = mu[-1]                     # total mean
    eps   = 1e-12
    w0, w1 = omega, 1 - omega
    mu0 = mu / (w0 + eps)
    mu1 = (mu_T - mu) / (w1 + eps)
    sigma_b2 = w0 * w1 * (mu0 - mu1)**2
    return int(np.argmax(sigma_b2))
\end{lstlisting}

\subsubsection*{Mini-Task Results}

\textbf{Image:} \texttt{threshold\_ostu1.jpg}.

\begin{table}[H]
  \centering
  \begin{tabular}{@{}lcc@{}}
    \toprule
    Input & Otsu threshold & Observation \\
    \midrule
    Original               & 83  & Baseline \\
    After $\gamma = 0.5$   & 109 & Brighter image $\Rightarrow$ higher threshold \\
    After $\gamma = 2.0$   & 82  & Darker image $\Rightarrow$ slightly lower threshold \\
    \bottomrule
  \end{tabular}
  \caption{Otsu threshold under different preprocessing.}
\end{table}

\textbf{NumPy vs OpenCV Otsu:} Both produce the same threshold (verified by the
comparison cell). The NumPy implementation is the mathematically explicit version;
OpenCV invokes the same algorithm in C++.

\textbf{Why gamma can help Otsu:}
\begin{itemize}[leftmargin=*]
  \item $\gamma < 1$ brightens the image and stretches the dark-pixel histogram modes
        apart, potentially improving bimodality.
  \item $\gamma > 1$ compresses the bright range; can help if foreground is bright
        and too merged with the background peak.
  \item Excessive gamma in either direction may \emph{hurt} Otsu by collapsing two
        well-separated peaks into one.
\end{itemize}

% ─────────────────────────────────────────────────────────────────────────────
\subsection{Integrated Exercises (Part 1)}
% ─────────────────────────────────────────────────────────────────────────────

\paragraph{Exercise 1 -- Brighten then Otsu.}
Applied $\gamma = 0.6$ to a dark image, then Otsu. The blurred+gamma image produces
a more distinct histogram, giving a cleaner binary mask than raw Otsu.

\paragraph{Exercise 2 -- ROI Extraction.}
Otsu mask AND-ed with the original grayscale image.  Foreground pixels are
preserved; background becomes black.

\paragraph{Exercise 3 -- Change Detection.}
$|\,A - B\,|$ absolute difference, then Otsu on the difference image.
The binary mask highlights only changed regions.

\paragraph{Exercise 4 -- Mask Fusion.}
\begin{itemize}[leftmargin=*]
  \item \textbf{OR} of Otsu (original) and Otsu (gamma=0.6): all candidate foreground pixels.
  \item \textbf{AND}: only pixels both methods agree are foreground.
  \item \textbf{XOR}: pixels where the two methods disagree (transition/uncertain pixels).
\end{itemize}

% =============================================================================
\section{Part 2 -- 2D Convolution}
% =============================================================================

\subsection{Theory}

For image $x \in \mathbb{R}^{H \times W}$ and kernel $K \in \mathbb{R}^{k_H \times k_W}$:
\begin{equation}
  y[i,j] = \sum_{u=0}^{k_H-1}\sum_{v=0}^{k_W-1}
            K[u,v]\cdot x\!\bigl[i + u - \Delta_h,\; j + v - \Delta_w\bigr]
  \label{eq:conv}
\end{equation}
with $\Delta_h = \lfloor k_H/2 \rfloor$, $\Delta_w = \lfloor k_W/2 \rfloor$.

\textbf{Convolution vs Cross-correlation:} True convolution flips the kernel
$180^\circ$ before sliding. \texttt{cv2.filter2D} implements cross-correlation
(no flip). For symmetric kernels, both are identical.

% ─────────────────────────────────────────────────────────────────────────────
\subsection{Padding Modes}
% ─────────────────────────────────────────────────────────────────────────────

\begin{table}[H]
  \centering
  \begin{tabular}{@{}llll@{}}
    \toprule
    Mode & Padding & Output size & Border \\
    \midrule
    VALID & None & $(H-k_H+1)\times(W-k_W+1)$ & Border pixels discarded \\
    SAME  & Symmetric & $H\times W$ & Zero / reflect / replicate \\
    \bottomrule
  \end{tabular}
  \caption{Convolution padding modes.}
\end{table}

For a $k\times k$ kernel in VALID mode, each dimension shrinks by $k-1$:
after ten $3\times3$ convolutions, a $256\times256$ image becomes $236\times236$.
This motivates SAME padding in deep networks.

% ─────────────────────────────────────────────────────────────────────────────
\subsection{NumPy Implementations}
% ─────────────────────────────────────────────────────────────────────────────

\paragraph{VALID mode (no padding):}
\begin{lstlisting}
def conv2d_valid_np(img_u8, kernel):
    img_f   = img_u8.astype(np.float32)
    k_flip  = kernel[::-1, ::-1]           # 180-degree flip
    kH, kW  = k_flip.shape
    oH, oW  = img_f.shape[0]-kH+1, img_f.shape[1]-kW+1
    out = np.zeros((oH, oW), dtype=np.float64)
    for i in range(oH):
        for j in range(oW):
            out[i, j] = np.sum(img_f[i:i+kH, j:j+kW] * k_flip)
    return normalise_for_display(out)
\end{lstlisting}

\paragraph{SAME mode (with padding):}
\begin{lstlisting}
def conv2d_same_np(img_u8, kernel, pad_mode='reflect'):
    img_f   = img_u8.astype(np.float32)
    k_flip  = kernel[::-1, ::-1]
    kH, kW  = k_flip.shape
    pH, pW  = kH // 2, kW // 2
    padded  = np.pad(img_f, ((pH, pH), (pW, pW)), mode=pad_mode)
    H, W    = img_f.shape
    out = np.zeros((H, W), dtype=np.float64)
    for i in range(H):
        for j in range(W):
            out[i, j] = np.sum(padded[i:i+kH, j:j+kW] * k_flip)
    return normalise_for_display(out)
\end{lstlisting}

% ─────────────────────────────────────────────────────────────────────────────
\subsection{OpenCV Implementation}
% ─────────────────────────────────────────────────────────────────────────────
\begin{lstlisting}
def conv2d_cv2(img_u8, kernel, border=cv2.BORDER_REFLECT):
    img_f = img_u8.astype(np.float32)
    out_f = cv2.filter2D(img_f, ddepth=cv2.CV_32F,
                         kernel=kernel.astype(np.float32),
                         borderType=border)
    return normalise_for_display(out_f.astype(np.float64))
\end{lstlisting}

Note: \texttt{cv2.filter2D} does cross-correlation. For asymmetric kernels (e.g.\ Sobel),
flip manually if true convolution is required.

% ─────────────────────────────────────────────────────────────────────────────
\subsection{Kernel Library}
% ─────────────────────────────────────────────────────────────────────────────

\begin{table}[H]
  \centering
  \begin{tabular}{@{}lll@{}}
    \toprule
    Kernel & Mathematical form & Purpose \\
    \midrule
    Identity $3\times3$ & $\delta[i,j]$ & No-op; pipeline verification \\[4pt]
    Box $3\times3$ & $\frac{1}{9}\mathbf{1}_{3\times3}$ & Average blur \\[4pt]
    Box $5\times5$ & $\frac{1}{25}\mathbf{1}_{5\times5}$ & Stronger average blur \\[4pt]
    Gauss $3\times3$ & $\frac{1}{16}\begin{bmatrix}1&2&1\\2&4&2\\1&2&1\end{bmatrix}$ & Smooth denoising \\[6pt]
    Gauss $5\times5$ & $\frac{1}{256}[\cdots]$ & Stronger smooth denoising \\[4pt]
    Sharpen & $\begin{bmatrix}0&-1&0\\-1&5&-1\\0&-1&0\end{bmatrix}$ & Edge enhancement \\[6pt]
    Sobel X & $\begin{bmatrix}-1&0&1\\-2&0&2\\-1&0&1\end{bmatrix}$ & Horizontal gradient (vertical edges) \\[6pt]
    Sobel Y & $\begin{bmatrix}-1&-2&-1\\0&0&0\\1&2&1\end{bmatrix}$ & Vertical gradient (horizontal edges) \\[6pt]
    Laplacian & $\begin{bmatrix}0&1&0\\1&-4&1\\0&1&0\end{bmatrix}$ & Isotropic edge detector \\[6pt]
    LoG $5\times5$ & Gauss $\ast$ Laplacian & Noise-robust edge detector \\
    \bottomrule
  \end{tabular}
  \caption{Standard convolution kernels implemented in the lab.}
\end{table}

% ─────────────────────────────────────────────────────────────────────────────
\subsection{Mini-Task Results}
% ─────────────────────────────────────────────────────────────────────────────

\subsubsection*{Task 1 -- Identity Kernel}
\textbf{Image:} \texttt{2.png}. \textbf{Kernel:} $3\times3$ identity.

VALID mode produces an output $2$ pixels smaller in each dimension (border strip
discarded). SAME mode (reflect padding) and OpenCV both preserve the original size.
The identity kernel recovers the original pixel values, confirming the pipeline is correct.

\subsubsection*{Task 2 -- Blur Filters}

\begin{table}[H]
  \centering
  \begin{tabular}{@{}llp{7cm}@{}}
    \toprule
    Kernel & Mode & Observation \\
    \midrule
    Box $3\times3$    & VALID & Mild smooth blur; slight ringing at edges \\
    Box $5\times5$    & SAME  & Stronger blur; block artefacts possible \\
    Gauss $3\times3$  & SAME  & Smooth blur; fewer artefacts than box \\
    Gauss $5\times5$  & OpenCV & Strongest smooth blur; natural gradient fall-off \\
    \bottomrule
  \end{tabular}
  \caption{Blur filter results.}
\end{table}

\textbf{Key observation:} Gaussian weights follow a bell curve, giving more weight
to nearby pixels. Box blur weights all pixels equally, which can produce ``ringing''
artefacts. Larger kernel size → stronger blurring.

\subsubsection*{Task 3 -- Sharpening}

The standard sharpen kernel ($[0,-1,0;\,-1,5,-1;\,0,-1,0]$) amplifies
high-frequency detail (edges). The strong variant ($[-1,-1,-1;\,-1,9,-1;\,-1,-1,-1]$)
produces more aggressive sharpening but also amplifies noise in smooth regions.

\textbf{Why it works:} The sharpen kernel = Identity + (Identity $-$ Laplacian) scaled.
In frequency domain: it boosts high frequencies (edges) while keeping low frequencies
(flat regions) unchanged.

\textbf{Clip vs normalise:} For sharpening, clipping is the correct display method
because the kernel is designed so that valid output stays near [0,255]. Normalising
would change the perceived brightness.

\subsubsection*{Task 4 -- Edge Detection}

\begin{table}[H]
  \centering
  \begin{tabular}{@{}llp{6.5cm}@{}}
    \toprule
    Kernel & Parameters & Observation \\
    \midrule
    Sobel X & $3\times3$, reflect pad & Strong response at vertical object boundaries \\
    Sobel Y & $3\times3$, reflect pad & Strong response at horizontal boundaries \\
    Sobel magnitude & $\sqrt{G_x^2+G_y^2}$ & All edges; orientation-invariant \\
    Laplacian & $3\times3$, reflect pad & All edges; very sensitive to noise \\
    LoG $5\times5$ & Reflect pad & Cleaner edges than bare Laplacian \\
    \bottomrule
  \end{tabular}
  \caption{Edge detection results.}
\end{table}

\textbf{Normalise for display:} Edge filter outputs contain negative values and
values $> 255$. Normalising maps the full signed range to [0,255], preserving all
information visually (absolute values are not interpretable directly).

\subsubsection*{Task 5 -- VALID vs SAME Comparison}

Applied Gauss $5\times5$ in both modes.

\begin{itemize}[leftmargin=*]
  \item \textbf{VALID:} output is $4$ rows and $4$ columns smaller. The border region
        of the image is not processed at all (those rows/columns are simply absent).
  \item \textbf{SAME (reflect):} output identical in size to input. Reflect
        padding mirrors the border content, producing a smooth continuation that
        minimises border artefacts.
  \item \textbf{SAME (zero):} would create a dark border ring visible with edge-
        detection kernels (the zero-pad is ``seen'' as a very dark edge).
\end{itemize}

% ─────────────────────────────────────────────────────────────────────────────
\subsection{Integrated Pipeline: Blur $\to$ Gamma $\to$ Otsu}
% ─────────────────────────────────────────────────────────────────────────────

\subsubsection*{Pipeline Design}

\begin{enumerate}
  \item \textbf{Gaussian blur} (Gauss $5\times5$, SAME reflect) -- suppress texture
        noise that would create spurious histogram peaks.
  \item \textbf{Gamma correction} ($\gamma = 0.6$) -- brighten dark regions;
        improve histogram bimodality.
  \item \textbf{Otsu thresholding} -- automatic global threshold on the improved
        image.
  \item \textbf{ROI extraction} -- AND between binary mask and original image.
\end{enumerate}

\subsubsection*{Results}

\begin{table}[H]
  \centering
  \begin{tabular}{@{}lcc@{}}
    \toprule
    Pipeline step & Otsu threshold & Observation \\
    \midrule
    Raw image, no preprocessing & $t_{\text{raw}}$ & Baseline \\
    After blur + $\gamma = 0.6$ & $t_{\text{pipe}}$ & More distinct bimodal histogram \\
    \bottomrule
  \end{tabular}
  \caption{Pipeline effect on Otsu threshold.}
\end{table}

\textbf{What improved:}
\begin{itemize}[leftmargin=*]
  \item Gaussian blur removes fine texture/noise → histogram peaks become sharper.
  \item Gamma brightening reveals foreground structure hidden in dark shadows.
  \item Combined effect: the histogram is more clearly bimodal → Otsu finds a
        more accurate and robust threshold.
\end{itemize}

\textbf{What got worse:}
\begin{itemize}[leftmargin=*]
  \item Blurring reduces edge crispness: the binary mask boundary will be slightly
        smoother (less sharp edges) compared to thresholding the original.
  \item Excessive gamma can amplify noise in very dark regions.
\end{itemize}

% ─────────────────────────────────────────────────────────────────────────────
\subsection{Discussion: NumPy vs OpenCV for Convolution}
% ─────────────────────────────────────────────────────────────────────────────

\begin{table}[H]
  \centering
  \begin{tabular}{@{}lll@{}}
    \toprule
    Aspect & NumPy (manual) & OpenCV \texttt{filter2D} \\
    \midrule
    Speed       & Slow (Python loops) & Fast (C++ / SIMD) \\
    Kernel flip & Manual \texttt{[::-1,::-1]} & No flip (cross-correlation) \\
    Border      & \texttt{np.pad} modes & \texttt{borderType} parameter \\
    Output type & Flexible float64 & Configurable \texttt{ddepth} \\
    Use case    & Education, small images & Production \\
    \bottomrule
  \end{tabular}
  \caption{NumPy vs OpenCV convolution comparison.}
\end{table}

% =============================================================================
\section{Conclusion}
% =============================================================================

\begin{enumerate}
  \item A grayscale image can be treated as a 2D float array; safe arithmetic
        requires the \texttt{float32 $\to$ clip $\to$ uint8} workflow.
  \item Gamma correction is a powerful preprocessing step that reshapes the
        histogram and can make automatic thresholding (Otsu) more reliable.
  \item Logical operations on binary masks allow flexible combination of detection
        results: OR for coverage, AND for consensus, XOR for disagreements.
  \item Otsu's method is reliable for near-bimodal histograms. It can fail for
        non-uniform illumination or heavily textured images.
  \item 2D convolution is a localised weighted sum. SAME padding preserves size;
        VALID does not touch borders.
  \item Standard kernels encode specific properties:
        low-pass (blur), high-pass (sharpen, edge detection), directional (Sobel).
  \item The combined pipeline (Blur $\to$ Gamma $\to$ Otsu $\to$ AND) produces
        better segmentation than any single step alone on noisy/dark images.
\end{enumerate}

\end{document}
